
\subsection{Alternativas de solución}
\label{sec:alternativas}

El análisis de alternativas tiene como propósito identificar enfoques tecnológicos para atender el problema formulado en la sección 1.1.2: ofrecer al comandante o coordinador de una emergencia información comprensible sobre la posible propagación de un incendio forestal, considerando las restricciones de tiempo, disponibilidad de datos y capacidad computacional.

Primero se examinan proyectos y sistemas existentes en Latinoamérica y otras regiones, identificando sus capacidades, beneficios para los usuarios y diferencias frente al proyecto. Posteriormente se plantean tres alternativas técnicas y se comparan con la solución seleccionada.

La comparación se realiza en función del alcance académico y territorial de este trabajo. Por tanto, no pretende demostrar que la propuesta sea superior a plataformas desarrolladas y utilizadas por organismos especializados.

\subsubsection{Proyectos y soluciones existentes}

\textbf{Cell2Fire+KITRAL en Chile.}

Cell2Fire es un entorno de simulación de propagación de incendios forestales desarrollado con participación de investigadores de la Universidad de Chile y colaboradores internacionales. Su adaptación Cell2Fire+KITRAL incorpora modelos de comportamiento del fuego orientados a condiciones forestales chilenas.

El sistema permite representar la evolución espacial de incendios considerando información territorial, combustible y condiciones ambientales. Una de sus características destacadas es su integración con modelos de optimización para estudiar estrategias de prevención, tratamientos de vegetación y localización de intervenciones \cite{vfpCell2FireChile}.

Para las entidades responsables de prevención, este enfoque permite comparar estrategias de manejo territorial antes de implementar medidas, contribuyendo a orientar recursos hacia sectores prioritarios.

Frente a este antecedente, el proyecto propuesto se concentra en un problema funcional diferente: consultar escenarios físicos FDS previamente calculados a partir de una situación observada. Su principal aporte esperado no es optimizar tratamientos preventivos, sino facilitar el acceso y la interpretación de resultados en un territorio específico de Bogotá.

\textbf{Proyecto de gemelos digitales multimodales del ISCI en Chile.}

En septiembre de 2026, el Instituto Sistemas Complejos de Ingeniería anunció la adjudicación de un proyecto FONDEF para desarrollar una plataforma inteligente de gestión de incendios forestales basada en gemelos digitales multimodales.

La iniciativa contempla integrar modelos de propagación, procesamiento paralelo mediante GPU y observaciones procedentes de drones y cámaras, con el propósito de actualizar escenarios predictivos y aportar información para la coordinación operativa y preventiva \cite{vfpISCI2026}.

El beneficio previsto para los organismos de emergencia consiste en mejorar la actualización de la información disponible y facilitar el análisis de escenarios conforme se reciben nuevas observaciones.

A diferencia de este proyecto, la propuesta del equipo no contempla construir un gemelo digital ni integrar directamente sensores, cámaras o drones. Se concentra en una biblioteca finita de resultados FDS y un mecanismo de consulta sobre condiciones representadas previamente.

La arquitectura precalculada limita la actualización física de los escenarios, pero también permite estudiar un flujo de software más acotado sin depender de infraestructura de observación permanente. El proyecto chileno debe considerarse una iniciativa en desarrollo y no un sistema operativo cuya validación ya esté demostrada.

\textbf{FIP Cerrado en Brasil.}

El proyecto FIP Cerrado, desarrollado con participación del Centro de Sensoriamento Remoto de la Universidade Federal de Minas Gerais, integra herramientas de monitoreo y simulación para apoyar la prevención y gestión de incendios en el bioma Cerrado.

Entre sus componentes se encuentra la modelación de propagación del fuego mediante información meteorológica, características de la vegetación y condiciones territoriales. La documentación institucional describe aproximaciones basadas en autómatas celulares para representar la evolución espacial del incendio \cite{vfpFIPCerrado}.

Para los responsables de áreas protegidas, estas herramientas contribuyen a reconocer zonas susceptibles, estudiar posibles recorridos del fuego y preparar intervenciones preventivas.

Su principal diferencia frente al proyecto es la metodología de propagación y la escala territorial de aplicación. La propuesta del equipo utilizará resultados FDS, con mayor énfasis en la preparación de escenarios físicos tridimensionales y su posterior consulta.

El antecedente brasileño demuestra la utilidad de adaptar la modelación al territorio, mientras que la propuesta busca estudiar una integración específica entre simulación física, catálogo geográfico y aplicación de escritorio.

\textbf{Sistema de Predicción de Peligro de Incendios Forestales de México.}

La Comisión Nacional Forestal de México desarrolla el Sistema de Predicción de Peligro de Incendios Forestales (SPPIF), que integra información ambiental, meteorológica, geográfica y satelital para producir pronósticos de peligro y apoyar el monitoreo de incendios \cite{vfpCONAFOR2026}.

El beneficio para los organismos de prevención y respuesta consiste en disponer de información territorial para identificar condiciones favorables a la ocurrencia de incendios, organizar vigilancia y orientar la preparación de recursos.

La plataforma aborda principalmente peligro y monitoreo, mientras que el proyecto del equipo pretende consultar resultados que representen la evolución temporal de un incendio bajo condiciones previamente simuladas.

Por ello, ambas aproximaciones son complementarias. Los mapas de peligro pueden contribuir a la preparación del territorio y a la selección de igniciones de referencia, pero no sustituyen directamente una biblioteca de escenarios de propagación.

\textbf{Wildfire Analyst de Technosylva.}

Technosylva desarrolla herramientas de modelación y análisis para incendios forestales. Entre sus capacidades se encuentran la ejecución de simulaciones de propagación, la representación cartográfica de resultados y el apoyo a la interpretación de escenarios durante emergencias.

La empresa documenta su colaboración con CAL FIRE en el desarrollo de herramientas orientadas al análisis de incendios y a la disponibilidad de información para decisiones operativas \cite{vfpTechnosylva}.

Para comandantes y organismos de respuesta, este tipo de plataforma ofrece la posibilidad de estudiar escenarios relacionados con condiciones específicas, interpretar posibles afectaciones y apoyar la coordinación de acciones.

Constituye uno de los referentes funcionales más cercanos a la propuesta. Sin embargo, la existencia de una plataforma operativa no garantiza su disponibilidad, licenciamiento o adaptación al conjunto de datos y requisitos de un prototipo académico en Bogotá.

El desarrollo propio permitirá controlar la estructura de la biblioteca FDS, los metadatos, las reglas de selección y la presentación de resultados para el sector elegido. Esta ventaja de control e investigación no implica que el prototipo ofrezca más funciones o mayor precisión que Wildfire Analyst.

\textbf{Wildland Fire Decision Support System en Estados Unidos.}

El Wildland Fire Decision Support System (WFDSS) integra información territorial, herramientas de análisis del comportamiento del fuego y documentación de decisiones para apoyar a responsables de la gestión de incendios en Estados Unidos \cite{vfpWFDSS}.

Su principal valor para los interesados radica en relacionar información científica con procedimientos de análisis y justificación de decisiones. No se limita a presentar una simulación, sino que forma parte de un entorno institucional de gestión.

Para el proyecto, WFDSS constituye un referente de diseño funcional y de comunicación de resultados. Sin embargo, adoptar directamente un sistema institucional de este tipo requeriría examinar sus condiciones de acceso, procesos y adecuación al contexto colombiano.

La solución propuesta tiene un alcance menor y se centra en un mecanismo específico de consulta de escenarios FDS, lo que permite estudiar su funcionamiento sin asumir la construcción de un sistema integral de gestión de emergencias.

\textbf{Spark Operational en Australia.}

El Commonwealth Scientific and Industrial Research Organisation (CSIRO) desarrolló Spark, un entorno de simulación de propagación de incendios que considera topografía, combustible y condiciones meteorológicas.

Spark Operational busca facilitar la generación de escenarios de propagación mediante métodos computacionales orientados a obtener resultados con rapidez, y se ha relacionado con necesidades de planificación y respuesta de organismos de emergencia \cite{vfpCSIROSpark}.

Para sus usuarios, la principal ventaja es la posibilidad de producir y examinar escenarios sin asumir necesariamente el mismo costo computacional de un modelo CFD tridimensional detallado.

Frente a Spark, la propuesta del equipo conserva el uso de FDS por su capacidad de representar determinados procesos físicos de forma explícita. El costo asociado se aborda mediante la generación anticipada de escenarios.

Esta decisión no demuestra que FDS sea más preciso o conveniente que Spark en cualquier situación. Su pertinencia depende de los fenómenos que se busquen representar, la disponibilidad de datos y la viabilidad de los cálculos para el sector seleccionado.

\textbf{B-FIREPREP en Portugal.}

El proyecto europeo B-FIREPREP se orientó al desarrollo de herramientas basadas en simulaciones estocásticas para anticipar situaciones peligrosas, identificar comunidades expuestas y apoyar la preparación frente a incendios forestales.

Su planteamiento contempla el análisis de múltiples escenarios y distintos horizontes temporales de anticipación, con énfasis en las necesidades de protección civil \cite{vfpBFIREPREP}.

Para los organismos responsables, el beneficio esperado consiste en identificar condiciones de riesgo y orientar la planificación de medidas preventivas y de respuesta.

Su relación con la propuesta se encuentra en la utilización de múltiples escenarios para apoyar decisiones. Sin embargo, B-FIREPREP se enfoca en anticipación y preparación territorial, mientras que el proyecto del equipo estudia la recuperación contextual de simulaciones FDS previamente ejecutadas.

La biblioteca propuesta permite controlar las condiciones de generación de cada resultado y su procedencia, aunque su cobertura estará limitada a las combinaciones preparadas y evaluadas.

En conjunto, estos antecedentes evidencian la existencia de diferentes aproximaciones a la simulación y gestión de incendios forestales. Sus beneficios dependen del usuario, la escala territorial, los modelos disponibles y las condiciones de aplicación.

La principal oportunidad del proyecto no consiste en reproducir íntegramente estas plataformas, sino en estudiar la viabilidad de una arquitectura concreta que transforme simulaciones físicas precalculadas en información consultable para un sector de los Cerros Orientales de Bogotá.

\subsubsection{Alternativas técnicas consideradas}

A partir de los antecedentes se identifican tres alternativas técnicas distintas de la propuesta seleccionada. Se consideran su costo de preparación, oportunidad de consulta, flexibilidad, dependencia de infraestructura y utilidad para los responsables de emergencias.

\textbf{Alternativa 1. Ejecución de simulaciones físicas bajo demanda.}

Esta alternativa consiste en desarrollar una aplicación que reciba una ubicación o perímetro observado junto con condiciones meteorológicas y ejecute una nueva simulación FDS para cada consulta.

Su principal ventaja sería permitir configurar el cálculo específicamente para la situación registrada, sin limitar la selección a escenarios previamente almacenados. Para el comandante o coordinador, esto podría ofrecer una representación ajustada a las condiciones de entrada disponibles.

Sin embargo, FDS resuelve numéricamente procesos físicos sobre una malla tridimensional. Su costo computacional depende del tamaño del dominio, la resolución espacial, la duración simulada y la complejidad de los modelos utilizados.

Como se expuso en la sección 5.1, incluso un dominio relativamente limitado puede contener millones de celdas. En consecuencia, ejecutar una simulación física completa durante una emergencia podría exigir recursos elevados y tiempos incompatibles con las necesidades inmediatas del usuario.

La simulación bajo demanda presenta, por tanto, un riesgo importante de inviabilidad operativa para el nivel de detalle y los recursos contemplados en el proyecto. No sería riguroso declarar que resulta imposible en todos los casos sin mediciones, pero tampoco puede asumirse que produciría resultados oportunos.

Frente a esta alternativa, la biblioteca precalculada desplaza la mayor parte del esfuerzo numérico a una etapa anterior a la consulta. Esta es su principal ventaja funcional: permite estudiar un mecanismo de respuesta basado en recuperación y procesamiento de resultados ya disponibles, aunque limitado por la cobertura de la biblioteca.

\textbf{Alternativa 2. Uso de modelos de propagación de menor costo computacional.}

La segunda alternativa consiste en desarrollar la aplicación alrededor de modelos de propagación simplificados o rápidos, como los empleados en FARSITE, Cell2Fire o Spark.

Su principal ventaja es la posibilidad de estudiar el avance del fuego en extensiones territoriales relativamente amplias sin resolver explícitamente todos los procesos de un modelo CFD tridimensional.

Para los organismos responsables, esta aproximación puede resultar útil cuando se requiere explorar grandes áreas, comparar escenarios de propagación o producir resultados mediante recursos computacionales más limitados.

Sin embargo, la representación física de estos modelos es diferente de FDS. Determinadas interacciones entre flujo de aire, combustión, geometría del entorno y combustible vegetal requieren aproximaciones que dependen del modelo seleccionado.

La propuesta del equipo es preferible para el objetivo específico de investigar una biblioteca basada en FDS, ya que conserva el motor físico elegido y permite estudiar su integración con un sistema de consulta.

Esta preferencia corresponde al alcance científico y académico definido, no a una superioridad general frente a los modelos rápidos. Estos últimos podrían ser más adecuados para otras escalas o restricciones de tiempo.

\textbf{Alternativa 3. Adaptación de una plataforma existente.}

Esta alternativa consiste en utilizar o adaptar una plataforma que ya integre simulaciones, información geográfica y herramientas de apoyo a decisiones, como Wildfire Analyst o WFDSS.

Para las entidades interesadas, una ventaja sería disponer de funciones previamente desarrolladas y, según el sistema, respaldadas por experiencias institucionales de utilización.

También podría reducir parte del esfuerzo de implementación si la plataforma permitiera incorporar los datos y requisitos del territorio.

Sin embargo, su utilización exigiría examinar licencias, condiciones de acceso, posibilidades de modificación, disponibilidad de funciones y compatibilidad con los modelos y conjuntos de datos requeridos.

Adicionalmente, adoptar una solución existente podría restringir el control del equipo sobre la generación de escenarios, la estructura de los resultados, la selección de casos y la experimentación con el procedimiento de agregación espacial.

La propuesta propia permite definir y comprobar estos componentes de forma explícita. Esto resulta especialmente pertinente porque el trabajo busca evaluar la viabilidad de un mecanismo de consulta sobre simulaciones FDS, no solamente disponer de una interfaz de visualización.

La principal ventaja frente a esta alternativa es, por tanto, la independencia de diseño y la trazabilidad del procedimiento experimental, aunque a cambio se asume un esfuerzo mayor de desarrollo e integración.

\subsubsection{Análisis comparativo de las alternativas}

La comparación se realiza mediante seis criterios: preparación computacional, tiempo de consulta, representación física, dependencia de infraestructura y datos, adaptación territorial y adecuación a los objetivos del proyecto.

Se utiliza una evaluación cualitativa porque todavía no existen mediciones experimentales equivalentes que permitan establecer diferencias cuantitativas de tiempo, costo o precisión entre las alternativas.

\begingroup
\small
\renewcommand{\arraystretch}{1.16}

\begin{longtable}{@{}p{0.18\linewidth}p{0.38\linewidth}p{0.35\linewidth}@{}}
\caption{Comparación cualitativa de alternativas de solución.}
\label{tab:comparacion-alternativas}\\
\hline
\textbf{Criterio} &
\textbf{Otras alternativas} &
\textbf{Biblioteca FDS precalculada} \\
\hline
\endfirsthead

\hline
\textbf{Criterio} &
\textbf{Otras alternativas} &
\textbf{Biblioteca FDS precalculada} \\
\hline
\endhead

Preparación computacional &
La simulación bajo demanda desplaza el cálculo al momento de consulta. Los modelos rápidos pueden reducirlo; las plataformas existentes dependen de su arquitectura. &
Requiere ejecutar y organizar previamente los escenarios, utilizando recursos HPC y almacenamiento. \\
\hline

Tiempo de consulta &
El cálculo FDS bajo demanda puede producir demoras elevadas. Los modelos rápidos y las plataformas existentes pueden ofrecer respuestas más oportunas, según su implementación. &
Recupera resultados almacenados sin ejecutar una nueva simulación FDS completa. La latencia efectiva debe medirse. \\
\hline

Representación física &
Los modelos rápidos utilizan simplificaciones diferentes de FDS. Las plataformas existentes pueden integrar diversos motores de simulación. &
Conserva resultados generados mediante FDS, sujetos a la calidad de sus entradas y configuración. \\
\hline

Infraestructura y datos &
FDS bajo demanda requiere recursos durante la emergencia. Los otros enfoques dependen de sus modelos, fuentes de información y condiciones de acceso. &
Concentra el uso intensivo de HPC en la preparación. Requiere mantener una biblioteca y datos territoriales compatibles. \\
\hline

Adaptación territorial &
Los modelos y plataformas deben parametrizarse para el sector. Una solución existente puede imponer restricciones técnicas o de acceso. &
Permite diseñar el catálogo y las reglas de consulta para el territorio estudiado, dentro de una cobertura finita. \\
\hline

Adecuación a objetivos &
Las alternativas ofrecen ventajas para simulación dinámica, análisis de grandes áreas o gestión institucional, pero no necesariamente permiten estudiar el mecanismo propuesto. &
Se ajusta al objetivo de desarrollar y evaluar una aplicación propia de consulta de escenarios FDS precalculados. \\
\hline
\end{longtable}
\endgroup

La comparación evidencia que no existe una solución universalmente superior. La conveniencia depende del propósito y las restricciones consideradas.

La simulación bajo demanda ofrece flexibilidad, pero introduce un costo computacional durante la consulta que podría comprometer su oportunidad. Los modelos rápidos pueden ser más adecuados para escalas territoriales amplias, mientras que las plataformas existentes pueden aportar funciones institucionales y experiencia operativa.

La biblioteca FDS precalculada ofrece una ventaja específica: separar la ejecución científica de la interacción del usuario. Esto permite desarrollar y evaluar un mecanismo de consulta sin exigir la repetición del cálculo físico completo durante cada solicitud.

Para los interesados en la aplicación, esta separación busca facilitar el acceso a escenarios comprensibles, manteniendo identificadas sus condiciones y limitaciones. Para el equipo académico, permite estudiar de forma controlada la generación, almacenamiento, recuperación, tratamiento espacial y visualización de resultados.

Sin embargo, su costo de preparación, cobertura territorial y validez de la aproximación deberán comprobarse experimentalmente antes de formular conclusiones de desempeño.

\subsubsection{Justificación de la alternativa seleccionada}

El equipo propone desarrollar una aplicación de escritorio basada en una biblioteca de escenarios FDS precalculados para un sector delimitado de los Cerros Orientales de Bogotá.

La aplicación permitirá que el comandante o coordinador ingrese una ubicación o un perímetro observado y utilice condiciones meteorológicas disponibles para recuperar simulaciones compatibles. Los resultados seleccionados serán procesados y presentados mediante una representación georreferenciada que facilite explorar posibles comportamientos del fuego.

La viabilidad técnica de esta propuesta se apoya en la separación entre la preparación de escenarios y su consulta. El cálculo físico se realizará antes del momento de utilización, aprovechando recursos HPC para generar resultados identificables y reproducibles. Posteriormente, la aplicación trabajará sobre información almacenada, evitando que el usuario tenga que configurar o ejecutar una simulación FDS completa cada vez que necesite consultar un escenario.

Esta arquitectura permite acotar el problema computacional y concentrar el desarrollo de software en funciones verificables: almacenamiento, catálogo, búsqueda, compatibilidad de condiciones, procesamiento espacial y visualización. Además, permite evaluar de manera independiente el costo de producción de la biblioteca y el tiempo de respuesta de la aplicación. Por estas razones, la solución presenta una ruta de implementación razonable dentro del alcance académico, aunque su viabilidad definitiva depende de los resultados de los pilotos.

En comparación con la simulación física bajo demanda, el principal beneficio esperado es disminuir la dependencia del procesamiento intensivo durante la consulta. Frente a modelos de propagación más simplificados, conserva el uso de FDS como motor físico seleccionado. Finalmente, frente a plataformas existentes, permite controlar el diseño y la evaluación del procedimiento específico de recuperación de escenarios.

Estas ventajas son pertinentes para los objetivos del proyecto, pero no implican que el prototipo alcance mayor precisión, oportunidad o madurez operativa que otras herramientas.

La solución también presenta limitaciones. Una biblioteca precalculada no puede representar todas las combinaciones posibles de ignición, combustible y condiciones ambientales. Su utilidad dependerá de la selección de casos, la representatividad de las entradas y la disponibilidad de escenarios compatibles con cada consulta.

La agregación de resultados desde un perímetro activo constituye una hipótesis de aproximación que deberá contrastarse mediante casos de referencia. Si no se demuestra una correspondencia defendible, su utilización tendrá que limitarse a las condiciones verificadas.

Estas restricciones no invalidan la propuesta como proyecto de grado. Por el contrario, permiten delimitar una contribución evaluable y establecer criterios para reconocer sus capacidades y límites. El alcance territorial reducido facilita concentrar la preparación y las comprobaciones en un conjunto de casos identificables, evitando extender conclusiones a escenarios no estudiados.

La aplicación será evaluada tanto por su funcionamiento como por la interpretación que los usuarios hagan de sus resultados. Una visualización comprensible y una consulta técnicamente correcta no equivalen, por sí mismas, a demostrar precisión científica ni autorización para su uso operativo.

En consecuencia, la alternativa seleccionada se considera pertinente porque permite investigar y desarrollar una integración concreta entre simulación física precalculada y apoyo informativo para la gestión de incendios. Su utilidad principal será facilitar la consulta de escenarios para el comandante o coordinador, con posibilidades complementarias de demostración y capacitación, sin sustituir el juicio profesional ni prometer resultados que todavía no hayan sido validados.
